% ===== PRÉAMBULE COMMUN — À COPIER TEL QUEL EN TÊTE DE CHAQUE FICHIER .tex =====
% Ne modifier QUE :
%   - les deux lignes \fancyhead[L] et \fancyhead[R]
%   - le bloc titre \begin{center}...\end{center} juste après \begin{document}
\documentclass[11pt,a4paper]{article}
\usepackage[utf8]{inputenc}
\usepackage[T1]{fontenc}
\usepackage{lmodern}
\usepackage[french]{babel}
\usepackage{amsmath,amssymb,mathtools}
\usepackage{icomma}
\usepackage{xcolor}
\usepackage{tikz}
\usepackage{pgfplots}
\usepackage[most]{tcolorbox}
\usepackage{enumitem}
\usepackage{array,booktabs,colortbl}
\usepackage[a4paper,margin=2.2cm,headheight=26pt,headsep=10pt]{geometry}
\usepackage{fancyhdr}
\usepackage[hidelinks]{hyperref}
\usetikzlibrary{arrows.meta,calc,angles,quotes,positioning,decorations.markings,intersections,backgrounds}
\pgfplotsset{compat=1.18}

\definecolor{primary}{RGB}{226,74,88}
\definecolor{primarydark}{RGB}{150,28,42}
\definecolor{methcol}{RGB}{40,90,150}
\definecolor{excol}{RGB}{0,135,90}
\definecolor{defcol}{RGB}{0,114,178}
\definecolor{attncol}{RGB}{200,40,55}
\definecolor{thmcol}{RGB}{0,140,120}
\definecolor{courbe}{RGB}{32,110,200}
\definecolor{courbeder}{RGB}{210,70,40}
\definecolor{tang}{RGB}{0,150,90}
\definecolor{grille}{RGB}{200,205,215}
\definecolor{plus}{RGB}{200,60,60}
\definecolor{moins}{RGB}{30,90,200}

\tcbset{boxbase/.style={enhanced,breakable,boxrule=0.9pt,arc=2.5pt,
  left=3mm,right=3mm,top=2.3mm,bottom=2.3mm,fonttitle=\bfseries,coltitle=white}}
\newtcolorbox{cle}[1][]{boxbase,colback=defcol!6,colframe=defcol,
  title={\textsc{À retenir}\ifx\relax#1\relax\relax\else\textmd{ : \itshape #1}\fi}}
\newtcolorbox{methode}[1][]{boxbase,colback=methcol!7,colframe=methcol,sharp corners,
  title={\textsc{Méthode}\ifx\relax#1\relax\relax\else\textmd{ --- \itshape #1}\fi}}
\newtcolorbox{exemplebox}[1][]{boxbase,colback=excol!5,colframe=excol,
  title={\textsc{Exemple}\ifx\relax#1\relax\relax\else\textmd{ : \itshape #1}\fi}}
\newtcolorbox{piege}[1][]{boxbase,colback=attncol!6,colframe=attncol,
  title={\textsc{Attention}\ifx\relax#1\relax\relax\else\textmd{ : \itshape #1}\fi}}
\newtcolorbox{exo}[1][]{boxbase,colback=primary!4,colframe=primary,
  title={\textsc{Exercice}\ifx\relax#1\relax\relax\else\textmd{~#1}\fi}}
\newtcolorbox{corr}[1][]{boxbase,colback=thmcol!5,colframe=thmcol,
  title={\textsc{Corrigé}\ifx\relax#1\relax\relax\else\textmd{~#1}\fi}}

\newcommand{\R}{\mathbb{R}}
\newcommand{\N}{\mathbb{N}}
\newcommand{\ds}{\displaystyle}
\newcommand{\tg}{\operatorname{tg}}
\newcommand{\cotg}{\operatorname{cotg}}
\newcommand{\arctg}{\operatorname{arctg}}
\newcommand{\arccotg}{\operatorname{arccotg}}
\newcommand{\limh}{\lim_{h\to 0}}

\pagestyle{fancy}\fancyhf{}
\fancyhead[L]{\small\textcolor{primarydark}{\textbf{Thème 2} — Formulaire de dérivation}}
\fancyhead[R]{\small\textcolor{primarydark}{\textsc{Corrigé --- Facile}}}
\fancyfoot[C]{\small\thepage}
\renewcommand{\headrulewidth}{0.8pt}

\begin{document}

\begin{center}
{\LARGE\bfseries\color{primarydark} Dérivées directes --- corrigé}\\[2pt]
{\color{primary}\rule{0.5\textwidth}{1.2pt}}\\[4pt]
{\small\itshape Niveau Facile --- constante, puissances, racines et inverses}
\end{center}
\vspace{2mm}

\begin{corr}[1 --- dérivée d'une constante]
$y=12$ est une constante, donc
\[
y'=0.
\]
\end{corr}

\begin{corr}[2 --- puissances entières]
Règle $(x^{n})'=n\,x^{n-1}$ et $(aU)'=aU'$ :
\[
y=x^{6}\ \Rightarrow\ y'=6\,x^{5};\qquad
y=4x^{3}\ \Rightarrow\ y'=4\cdot 3\,x^{2}=12\,x^{2}.
\]
\end{corr}

\begin{corr}[3 --- combinaison de puissances]
On dérive terme à terme (règle de la somme et du multiple) :
\[
y=5x^{4}-3x^{2}+7x-9
\ \Rightarrow\
y'=5\cdot4x^{3}-3\cdot2x+7\cdot1-0=20x^{3}-6x+7.
\]
\end{corr}

\begin{corr}[4 --- racines réécrites en puissances]
\textbf{a)} $y=\sqrt{x}=x^{1/2}$, donc
\[
y'=\tfrac12\,x^{-1/2}=\frac{1}{2\,x^{1/2}}=\frac{1}{2\sqrt{x}}.
\]
\textbf{b)} $y=\sqrt[3]{x}=x^{1/3}$, donc
\[
y'=\tfrac13\,x^{-2/3}=\frac{1}{3\,x^{2/3}}=\frac{1}{3\sqrt[3]{x^{2}}}.
\]
\end{corr}

\begin{corr}[5 --- inverses réécrits en puissances]
\textbf{a)} $y=\dfrac{1}{x^{3}}=x^{-3}$, donc
\[
y'=-3\,x^{-4}=-\frac{3}{x^{4}}.
\]
\textbf{b)} $y=\dfrac{2}{x}=2x^{-1}$, donc
\[
y'=2\cdot(-1)\,x^{-2}=-2\,x^{-2}=-\frac{2}{x^{2}}.
\]
\end{corr}

\end{document}
